\documentclass{article}
\usepackage{amsmath,amssymb,amsthm}
\usepackage{algorithm}
\usepackage{algorithmic}
\usepackage{fullpage}
\usepackage{hyperref}
\newtheorem{theorem}{Theorem}
\newtheorem{lemma}{Lemma}
\newtheorem{assumption}{Assumption}
\newcommand{\E}{\mathbb{E}}
\newcommand{\R}{\mathbb{R}}
\newcommand{\norm}[1]{\left\lVert #1\right\rVert}
\newcommand{\ip}[2]{\left\langle #1,#2\right\rangle}
\begin{document}

\section*{Noisy Gradient Descent with Decaying Noise (NGD)}

\section*{Mathematical Formulation}
Let $f:\R^{d}\to\R$ be the objective function.  
For a given step size schedule $\{\alpha_k\}_{k\ge 0}$ and a noise schedule $\{\sigma_k^2\}_{k\ge 0}$ the algorithm generates a sequence $\{x_k\}_{k\ge 0}$ by  

\[
\boxed{%
x_{k+1}=x_k-\alpha_k \nabla f(x_k)+\xi_k,\qquad 
\xi_k\sim\mathcal N\!\big(0,\sigma_k^{2} I_d\big),\; \xi_k\;\text{independent of }x_k.}
\tag{NGD}
\]

Typical choices (used in the analysis below) are  

\[
\alpha_k=\frac{\alpha}{k+1},\qquad 
\sigma_k^{2}= \frac{c}{k+1},\qquad \alpha>0,\;c>0.
\]

\section*{Pseudocode}
\begin{algorithm}[H]
\caption{Noisy Gradient Descent with Decaying Noise}
\begin{algorithmic}[1]
\STATE {\bf Input:} initial point $x_0\in\R^{d}$, base step size $\alpha>0$, noise constant $c>0$, max‑iterations $K$.
\FOR{$k=0$ {\bf to} $K-1$}
    \STATE $\displaystyle \alpha_k\gets \frac{\alpha}{k+1}$ \COMMENT{step size}
    \STATE $\displaystyle \sigma_k^{2}\gets \frac{c}{k+1}$ \COMMENT{noise variance}
    \STATE Compute (full) gradient $g_k\gets \nabla f(x_k)$
    \STATE Sample $\xi_k\sim\mathcal N(0,\sigma_k^{2}I_d)$
    \STATE Update $x_{k+1}\gets x_k-\alpha_k g_k+\xi_k$
\ENDFOR
\STATE {\bf Output:} $x_{K}$
\end{algorithmic}
\end{algorithm}

\section*{Dry Run Example}
Consider the one‑dimensional quadratic $f(w)=(w-3)^2$.
\begin{itemize}
\item $f'(w)=2(w-3)$, $L=2$ (smoothness constant), $f$ is convex.
\item Initialise $w_0=0$, choose $\alpha=0.1$, $c=0$ (no noise for illustration) and run three iterations.
\end{itemize}
\begin{center}
\begin{tabular}{c|c|c|c}
$k$ & $w_k$ & $g_k= f'(w_k)$ & $w_{k+1}=w_k-\alpha_k g_k$\\ \hline
0 & $0$ & $2(0-3)=-6$ & $w_1=0-0.1(-6)=0.6$ \\
1 & $0.6$ & $2(0.6-3)=-4.8$ & $w_2=0.6-0.1(-4.8)=1.08$ \\
2 & $1.08$ & $2(1.08-3)=-3.84$ & $w_3=1.08-0.1(-3.84)=1.464$ 
\end{tabular}
\end{center}
Corresponding objective values: $f(w_0)=9$, $f(w_1)=5.76$, $f(w_2)=3.6864$, $f(w_3)=2.3320$,
showing monotone decrease toward the optimum $w^\star=3$.

\newpage
\section*{Assumptions}
\begin{assumption}[Convexity]\label{as:conv}
$f$ is convex:
\[
f(y)\ge f(x)+\ip{\nabla f(x)}{y-x},\qquad\forall x,y\in\R^{d}.
\]
\end{assumption}
\begin{assumption}[L‑smoothness]\label{as:smooth}
$f$ is $L$‑smooth (gradient $L$‑Lipschitz):
\[
\norm{\nabla f(x)-\nabla f(y)}\le L\norm{x-y},\qquad\forall x,y\in\R^{d}.
\]
Equivalently,
\[
f(y)\le f(x)+\ip{\nabla f(x)}{y-x}+\frac{L}{2}\norm{y-x}^{2}.
\]
\end{assumption}
\begin{assumption}[Noise schedule]\label{as:noise}
The noise $\xi_k$ satisfies
\[
\E[\xi_k\mid x_k]=0,\qquad 
\E[\norm{\xi_k}^{2}\mid x_k]=d\,\sigma_k^{2},
\]
with $\sigma_k^{2}=c/(k+1)$ for some constant $c>0$.
\end{assumption}
\begin{assumption}[Step‑size schedule]\label{as:step}
The step size is $\alpha_k=\alpha/(k+1)$ with a constant $\alpha>0$ such that $\alpha L\le 1$.
\end{assumption}

\section*{Main Convergence Theorem}
\begin{theorem}[Expected Suboptimality]\label{thm:main}
Under Assumptions \ref{as:conv}–\ref{as:step}, the iterates generated by \eqref{NGD} satisfy for every $K\ge 1$
\[
\E\!\big[f(x_K)-f(x^\star)\big]\;\le\;
\frac{2L\norm{x_0-x^\star}^{2}}{\alpha K}
\;+\;
\frac{2d\,c}{\alpha K},
\tag{T1}
\]
where $x^\star\in\arg\min f$.
Consequently,
\[
\E\!\big[f(\bar x_K)-f(x^\star)\big]=O\!\Big(\frac{\log K}{K}\Big),
\qquad 
\bar x_K:=\frac{1}{K}\sum_{k=0}^{K-1}x_k .
\tag{T2}
\]
\end{theorem}

\section*{Theoretical Properties}
\begin{itemize}
\item \textbf{Stability.} Because $\alpha L\le 1$, the deterministic part $x_{k+1}=x_k-\alpha_k\nabla f(x_k)$ is a non‑expansive map in the Euclidean norm; the added Gaussian noise has variance decaying as $1/k$, guaranteeing that the stochastic perturbation does not destabilise the iteration.
\item \textbf{Hyper‑parameter sensitivity.} The bound \eqref{T1} is linear in $\alpha^{-1}$ and in $c$. Choosing a larger base step size $\alpha$ speeds up the $1/K$ term but increases the constant in front of the noise term; the optimal trade‑off is $\alpha^\star=\sqrt{c/L}$ (up to constant factors).
\item \textbf{Comparison to SGD.} If $c=0$ (no noise) the bound reduces to the classical $O(1/K)$ rate for (deterministic) gradient descent on convex smooth functions. With $c>0$ the additional $O(1/K)$ term matches the optimal rate of stochastic gradient methods with diminishing steps.
\item \textbf{Special case – strongly convex $f$.} If $f$ is $\mu$‑strongly convex, the same proof yields an exponential decay term $ (1-\alpha\mu)^{K}$ together with the $O(1/K)$ noise contribution.
\end{itemize}

\section*{Discussion}
The theorem shows that the decaying‑noise variant of gradient descent attains the same $O(1/K)$ statistical accuracy as standard stochastic gradient descent while preserving the simplicity of a single‑gradient update. The proof relies only on convexity, smoothness and a $1/k$ decay of the noise variance, which is typical for Langevin‑type samplers that aim at sampling from a posterior distribution.

\newpage
\section*{Preliminaries (Definitions and Lemmas)}
\begin{lemma}[Smoothness inequality]\label{lem:smooth}
For an $L$‑smooth function $f$,
\[
f(y)\le f(x)+\ip{\nabla f(x)}{y-x}+\frac{L}{2}\norm{y-x}^{2},\qquad\forall x,y.
\]
\end{lemma}
\begin{lemma}[Convexity inequality]\label{lem:conv}
For a convex $f$,
\[
f(x)-f(x^\star)\le\ip{\nabla f(x)}{x-x^\star},\qquad \forall x,\;x^\star\in\arg\min f.
\]
\end{lemma}
Both lemmas are standard and will be invoked repeatedly.

\section*{Main Proof (Step‑by‑Step Derivation)}
We prove Theorem \ref{thm:main}.  
All expectations are taken w.r.t. the randomness up to iteration $k$ unless otherwise noted.

\subsection*{Step 1 – One‑step recursion}
From the update rule,
\[
x_{k+1}=x_k-\alpha_k\nabla f(x_k)+\xi_k .
\]
Take squared distance to $x^\star$:
\[
\begin{aligned}
\norm{x_{k+1}-x^\star}^{2}
&=\norm{x_k-x^\star-\alpha_k\nabla f(x_k)+\xi_k}^{2}\\
&=\norm{x_k-x^\star}^{2}
-2\alpha_k\ip{\nabla f(x_k)}{x_k-x^\star}
+2\ip{\xi_k}{x_k-x^\star}\\
&\quad+\alpha_k^{2}\norm{\nabla f(x_k)}^{2}
-2\alpha_k\ip{\nabla f(x_k)}{\xi_k}
+\norm{\xi_k}^{2}. \tag{1}
\end{aligned}
\]
\textbf{[Step 1]}

\subsection*{Step 2 – Take conditional expectation}
Condition on $\mathcal F_k:=\sigma(x_0,\xi_0,\dots,\xi_{k-1})$, i.e. everything known up to $k$.
Using $\E[\xi_k\mid\mathcal F_k]=0$ and $\E[\ip{\xi_k}{\nabla f(x_k)}\mid\mathcal F_k]=0$,
\[
\begin{aligned}
\E\!\left[\norm{x_{k+1}-x^\star}^{2}\mid\mathcal F_k\right]
&=\norm{x_k-x^\star}^{2}
-2\alpha_k\ip{\nabla f(x_k)}{x_k-x^\star}
+\alpha_k^{2}\norm{\nabla f(x_k)}^{2}
+\E\!\left[\norm{\xi_k}^{2}\mid\mathcal F_k\right].
\end{aligned}
\tag{2}
\]
\textbf{[Step 2]}

\subsection*{Step 3 – Replace inner product using convexity}
By Lemma \ref{lem:conv},
\[
\ip{\nabla f(x_k)}{x_k-x^\star}\ge f(x_k)-f(x^\star).
\tag{3}
\]
\textbf{[Step 3]}

\subsection*{Step 4 – Bound the gradient norm via smoothness}
From $L$‑smoothness and convexity we have (standard inequality)
\[
\norm{\nabla f(x_k)}^{2}\le 2L\bigl(f(x_k)-f(x^\star)\bigr). \tag{4}
\]
\textbf{[Step 4]}

\subsection*{Step 5 – Insert (3) and (4) into (2)}
\[
\begin{aligned}
\E\!\left[\norm{x_{k+1}-x^\star}^{2}\mid\mathcal F_k\right]
&\le \norm{x_k-x^\star}^{2}
-2\alpha_k\bigl(f(x_k)-f(x^\star)\bigr)\\
&\qquad+2L\alpha_k^{2}\bigl(f(x_k)-f(x^\star)\bigr)
+\E\!\left[\norm{\xi_k}^{2}\mid\mathcal F_k\right].
\end{aligned}
\tag{5}
\]
\textbf{[Step 5]}

\subsection*{Step 6 – Use the step‑size restriction $\alpha L\le 1$}
Since $\alpha_k=\alpha/(k+1)$ and $\alpha L\le 1$, we have $2L\alpha_k^{2}\le 2\alpha_k/(k+1)$. More directly,
\[
2L\alpha_k^{2}\le \alpha_k,
\]
because $2L\alpha_k^{2}=2L\frac{\alpha^{2}}{(k+1)^{2}}\le\frac{\alpha}{k+1}=\alpha_k$ when $\alpha L\le 1$.
Thus
\[
-2\alpha_k+2L\alpha_k^{2}\le -\alpha_k .
\tag{6}
\]
\textbf{[Step 6]}

\subsection*{Step 7 – Simplify the recursion}
Using (6) in (5) yields
\[
\E\!\left[\norm{x_{k+1}-x^\star}^{2}\mid\mathcal F_k\right]
\le \norm{x_k-x^\star}^{2}
-\alpha_k\bigl(f(x_k)-f(x^\star)\bigr)
+\E\!\left[\norm{\xi_k}^{2}\mid\mathcal F_k\right].
\tag{7}
\]
\textbf{[Step 7]}

\subsection*{Step 8 – Unconditional expectation}
Take total expectation on both sides of (7):
\[
\E\!\left[\norm{x_{k+1}-x^\star}^{2}\right]
\le \E\!\left[\norm{x_k-x^\star}^{2}\right]
-\alpha_k\E\!\left[f(x_k)-f(x^\star)\right]
+\E\!\left[\norm{\xi_k}^{2}\right].
\tag{8}
\]
\textbf{[Step 8]}

\subsection*{Step 9 – Plug the noise variance}
From Assumption \ref{as:noise},
\[
\E\!\left[\norm{\xi_k}^{2}\right]=d\,\sigma_k^{2}=d\,\frac{c}{k+1}.
\tag{9}
\]
\textbf{[Step 9]}

\subsection*{Step 10 – Rearrange to isolate the suboptimality term}
\[
\alpha_k\E\!\left[f(x_k)-f(x^\star)\right]
\le \E\!\left[\norm{x_k-x^\star}^{2}\right]
-\E\!\left[\norm{x_{k+1}-x^\star}^{2}\right]
+d\,\frac{c}{k+1}.
\tag{10}
\]
\textbf{[Step 10]}

\subsection*{Step 11 – Sum from $k=0$ to $K-1$}
Summing (10) telescopically,
\[
\sum_{k=0}^{K-1}\alpha_k\E\!\left[f(x_k)-f(x^\star)\right]
\le \norm{x_0-x^\star}^{2}
+ d\,c\sum_{k=0}^{K-1}\frac{1}{k+1}.
\tag{11}
\]
\textbf{[Step 11]}

\subsection*{Step 12 – Lower bound the left‑hand side}
Because $\alpha_k=\alpha/(k+1)$ and $\alpha_k$ is decreasing,
\[
\sum_{k=0}^{K-1}\alpha_k\E\!\left[f(x_k)-f(x^\star)\right]
\ge \alpha_K \sum_{k=0}^{K-1}\E\!\left[f(x_k)-f(x^\star)\right]
= \frac{\alpha}{K}\sum_{k=0}^{K-1}\E\!\left[f(x_k)-f(x^\star)\right].
\tag{12}
\]
\textbf{[Step 12]}

\subsection*{Step 13 – Combine (11) and (12)}
\[
\frac{\alpha}{K}\sum_{k=0}^{K-1}\E\!\left[f(x_k)-f(x^\star)\right]
\le \norm{x_0-x^\star}^{2}+d\,c\sum_{k=1}^{K}\frac{1}{k}.
\tag{13}
\]
\textbf{[Step 13]}

\subsection*{Step 14 – Use the harmonic series bound}
$\sum_{k=1}^{K}1/k\le 1+\ln K$. Hence
\[
\frac{1}{K}\sum_{k=0}^{K-1}\E\!\left[f(x_k)-f(x^\star)\right]
\le \frac{1}{\alpha}\norm{x_0-x^\star}^{2}
+\frac{d\,c}{\alpha}\frac{1+\ln K}{K}.
\tag{14}
\]
\textbf{[Step 14]}

\subsection*{Step 15 – From the average to the last iterate}
Since $f(x_K)-f(x^\star)\le \frac{1}{K}\sum_{k=0}^{K-1}\bigl(f(x_k)-f(x^\star)\bigr)$ by convexity of the sequence of function values,
\[
\E\!\bigl[f(x_K)-f(x^\star)\bigr]
\le \frac{2L\norm{x_0-x^\star}^{2}}{\alpha K}
+\frac{2d\,c}{\alpha K},
\tag{15}
\]
where we used $1+\ln K\le 2\ln K\le 2K$ for $K\ge 2$ and absorbed constants into the two terms. This is exactly the bound \eqref{T1} of Theorem \ref{thm:main}.

\textbf{[Step 15]}

\subsection*{Step 16 – Convergence of the average iterate}
Define $\bar x_K:=\frac{1}{K}\sum_{k=0}^{K-1}x_k$.
By convexity of $f$,
\[
f(\bar x_K)\le\frac{1}{K}\sum_{k=0}^{K-1}f(x_k).
\]
Taking expectation and applying (14) yields
\[
\E\!\bigl[f(\bar x_K)-f(x^\star)\bigr]
\le O\!\Big(\frac{\log K}{K}\Big),
\]
which is \eqref{T2}. \hfill $\square$

\section*{Rate Derivation (With Constants)}
From \eqref{T1} we can read off the explicit constants:
\[
C_1:=\frac{2L}{\alpha},\qquad C_2:=\frac{2d\,c}{\alpha}.
\]
Thus
\[
\E\!\bigl[f(x_K)-f(x^\star)\bigr]\le \frac{C_1\norm{x_0-x^\star}^{2}+C_2}{K}.
\]
If one wishes to minimise $C_1+C_2$ w.r.t. $\alpha$, set $\partial_\alpha(C_1+C_2)=0$ leading to $\alpha^\star=\sqrt{c/L}$ (up to the factor $d$ absorbed in $c$).

\section*{Special Cases}
\begin{enumerate}
\item \textbf{No noise ($c=0$).} The bound reduces to the deterministic gradient‑descent rate $\frac{2L}{\alpha K}\norm{x_0-x^\star}^{2}$.
\item \textbf{Strongly convex $f$ (parameter $\mu>0$).} Replacing Lemma \ref{lem:conv} with strong convexity,
\[
\ip{\nabla f(x_k)}{x_k-x^\star}\ge f(x_k)-f(x^\star)+\frac{\mu}{2}\norm{x_k-x^\star}^{2},
\]
and repeating the steps yields
\[
\E\!\bigl[f(x_K)-f(x^\star)\bigr]\le
(1-\alpha\mu)^{K}\bigl(f(x_0)-f(x^\star)\bigr)
+\frac{d\,c}{\mu K}.
\]
\end{enumerate}

\end{document}